\section{Main ideas}\label{main-ideas}

We aim to design a type system that allows its users to
reuse native syntax of a system for special kinds of
user-defined ``domain-specific languages''. This
approach has several advantages:

\begin{itemize}
    \item From the pragmatic standpoint, it would allow
        DSL designers to use a syntax which is already
        familiar to their users;
    \item From type-theoretical standpoint, type
        theorists could internalize metatheoretical
        properties of such type system directly inside
        itself;
    \item From philosophical standpoint, it would be
        particularly satisfying to have a type system
        capable of ``eating itself''.
\end{itemize}

To define a new DSL in such system, it is enough to
specify how to interpret the native syntax inside your
domain of interest. It is natural to expect that
provided interpretations respect definitional equalities
of a system; finally, we want interpretations to be
definable inside a system, so we call them
\textit{internal models} of a system, or just
\textit{model} for short.

Now, before actually specifying a system, let's find out
what we cannot hope to achieve.

\subsection{No-go theorem of internal interpretation}

We're going to work inside a classical type-theoretical
framework which introduces a sort of contexts $\Gamma$,
terms $t,t'$ and types $T$ linked by the typing relation
$\Gamma\vdash t:T$; computation is specified by one-step
reduction relation $t\to t'$ on terms. In this section
we'll show that, if we want a practical system with
internal interpretation as discussed in the start of the
section \ref{main-ideas}, we have to either a)
reinterpret closed terms only; b) forbid reductions
inside reintrerpreted terms altogether; c) severely
restrict a class of definable models.

\begin{definition}
    [Na\"{\i}ve computational reinterpretation]
    We call a type operator $F[\_]$ a na\"{\i}ve
    computational reinterpretation of types if it
    respects computation on types (that is, if there is
    a reduction relation on types $(\to)$ and $T\to T'$,
    then $F[T]\to^* F[T']$), and computes as a primitive
    recursive function on normal forms of types.
    Analogously, term with a hole
    $\llbracket\_\rrbracket$ is a na\"{\i}ve
    computational reinterpretation of terms if it
    respects computation on terms, computes as a
    primitive recursive function on normal forms of
    terms and, if $\Gamma\vdash t:T$, then
    $\Gamma\vdash\llbracket t\rrbracket:F[T]$.
\end{definition}

\begin{theorem}[No-go theorem of naive reinterpretation]
\label{no-go}
    In a confluent type system with a type of functions
    and na\"{\i}ve computational reinterpretation, the
    following should hold for arbitrary well-typed $s$
    and $t$:
    \[
        \mathtt{app}\llbracket\lambda x.t\rrbracket
        \llbracket s \rrbracket \equiv
        \llbracket t[x\coloneqq s] \rrbracket
    \]
    where $\equiv$ is definitional equality and
    $\mathtt{app}$ is how $\llbracket\_\rrbracket$
    acts on applications.
\end{theorem}

Principally, it means that \textbf{we cannot use
propositional equality to enforce equational properties
of models}. The proof is by direct example:

\begin{proof}
    We are given well-typed $\Gamma\vdash s:S$ and
    $\Gamma,x:S\vdash t:T$. Consider the following term:
    \[
        \Gamma,x:S\to T,y:S\vdash
            \llbracket x\,y\rrbracket:F[T]
    \]
    Term under interpretation is in normal form, hence
    the whole term computes to $\mathtt{app}\llbracket x
    \rrbracket\llbracket y\rrbracket$.
    Substituting $\lambda x.t$ for $x$ and $s$ for $y$,
    we get $\mathtt{app}\llbracket\lambda x.t\rrbracket
    \llbracket s \rrbracket$.

    On the other hand, if we perform substitution first,
    we would get $\Gamma\vdash\llbracket(\lambda x.t)s
    \rrbracket:F[T]$. Since na\"{\i}ve reinterpretation
    respects reduction, this reduces to
    $\llbracket t[x\coloneqq s] \rrbracket$.

    In a confluent type system with functions,
    multi-step reduction commutes with substitution,
    therefore two terms are definitionally equivalent.
\end{proof}

Fundamentally, the proof above exploits openness of the
original term; if we restrict reinterpretation to
compute on closed terms only, we can arrive at a sound
system without requiring definitional equalities in
models, as is exemplified by our non-dependently-typed
system STLC-R described in section \ref{STLC-R}.
However, such a tradeoff is unacceptable in a system
with dependent types which has to compute on open terms
in order to have a flexible typing discipline. This
leaves us with three options:

\begin{enumerate}
    \item Introduce a ``modality of global elements''
        $\Box$ which reifies a metatheoretical notion of
        a closed term in a dependent type system, then
        allow to reinterpret boxed terms only.
    \item Forbid computations under
        $\llbracket\_\rrbracket$ unless term is closed.
        To allow reasoning on open terms, equations
        describing behavior of interpretation are
        exposed as axioms in equality type.
    \item Restrict models to those where equalities hold
        definitionally rather than propositionally.
\end{enumerate}

While the first approach is most desirable, it
unfortunately still doesn't work, as is shown in section
\ref{global-elements}. The second one looks consistent
with propositional definition of models and is left for
future work; in \ref{MLTT-R}, we explore the third
approach as it seems easier to work with in practice
thanks to more flexible computation rules.
